\section{Momentum theory over an actuator disk}
\label{sec:momentum}

Everything downstream rests on one result: the power required to hold a mass in
the air, as a function of how much air the machine can reach. We derive it from
a control volume rather than quote it, because the exponents that appear are
the entire argument of \cref{sec:fold} and it matters that they are earned.

\subsection{The model}

\begin{model}[Rankine-Froude actuator disk]
\label{mod:disk}
The rotor is replaced by a permeable disk of area $A$ across which the static
pressure jumps by $\Delta p$ while the velocity is continuous. The flow is
axisymmetric, the fluid is at rest far upstream, and the slipstream is a
well-defined streamtube. Rotation imparted to the wake, viscous losses and the
finite number of blades are all excluded; they return in \cref{sec:blade}.
\end{model}

This is the classical model of Rankine and W.~Froude, given in the form used
here by Glauert~\cite{glauert1935} and, in the rotorcraft setting, by
Leishman~\cite[Ch.~2]{leishman2006} and Johnson~\cite[Ch.~2]{johnson1980}.

Let $\vind$ denote the axial velocity induced \emph{at the disk} and $w$ the
axial velocity in the far wake, both measured relative to the undisturbed
fluid. Let $A_\infty$ be the cross-sectional area of the far wake.
\Cref{fig:streamtube} shows the control volume and the four stations used
below.

\begin{figure}[tbp]
\centering
% Actuator-disk streamtube in hover, flow downward. The boundary radius is
% drawn from continuity, r(x)^2 u(x) = R^2 v_i, with the axial velocity of
% the vortex-cylinder solution on the axis, u/v_i = 1 + xi/sqrt(1+xi^2),
% xi = x/R positive downstream. One-dimensional, but the shape is earned.
\begin{tikzpicture}[scale=1.25,
  declare function={u(\s) = 1 + \s/sqrt(1+\s*\s);
                    trad(\s) = 1/sqrt(u(\s));}]
  \def\xa{-1.05} \def\xb{3.6}
  % streamtube boundaries (y = -xi, flow downward)
  \fill[cA!6] plot[domain=\xa:\xb, samples=80] ({ trad(\x)}, {-\x})
      -- plot[domain=\xb:\xa, samples=80] ({-trad(\x)}, {-\x}) -- cycle;
  \draw[cA, line width=0.9pt] plot[domain=\xa:\xb, samples=80] ({ trad(\x)}, {-\x});
  \draw[cA, line width=0.9pt] plot[domain=\xa:\xb, samples=80] ({-trad(\x)}, {-\x});
  % interior streamlines, scaled boundaries
  \foreach \f in {0.35, 0.7}{
    \draw[cA!55, line width=0.4pt] plot[domain=\xa:\xb, samples=60] ({ \f*trad(\x)}, {-\x});
    \draw[cA!55, line width=0.4pt] plot[domain=\xa:\xb, samples=60] ({-\f*trad(\x)}, {-\x});
  }
  % axis
  \draw[cG, dash dot, line width=0.4pt] (0,0.95) -- (0,-3.75);
  % actuator disk
  \draw[line width=2.4pt, black] (-1,0) -- (1,0);
  \node[lbl, anchor=west] at (1.28,0.1) {disk, area $A$};
  \draw[vec, cD, line width=1.2pt] (0,0.05) -- (0,0.72) node[right, lbl] {$T$};
  % control volume
  \draw[construct] (-2.45,1.0) rectangle (2.45,-3.72);
  \node[lbl, cG, anchor=north west] at (-2.42,-3.3) {control volume};
  % stations
  \foreach \y/\n in {0.92/0, 0.13/1, -0.13/2, -3.55/3}
    \node[lbl, anchor=east] at (-2.55,\y) {\n};
  \draw[cG, line width=0.3pt] (-2.45,0.13) -- (-1.05,0.13);
  \draw[cG, line width=0.3pt] (-2.45,-0.13) -- (-1.05,-0.13);
  % velocity arrows: 0 far upstream, v_i at the disk, w = 2 v_i far downstream
  \node[lbl, anchor=west, align=left] at (3.3,0.92)
      {0: $p_\infty$, at rest};
  \node[lbl, anchor=west, align=left] at (3.3,0.28)
      {1: $p_1 = p_\infty - \tfrac12\rho\vind^2$, speed $\vind$};
  \node[lbl, anchor=west, align=left] at (3.3,-0.28)
      {2: $p_2 = p_1 + \Delta p$, speed $\vind$};
  \node[lbl, anchor=west, align=left] at (3.3,-3.55)
      {3: $p_\infty$, speed $w = 2\vind$, area $A_\infty = A/2$};
  \foreach \xx in {-0.55, 0, 0.55}
    \draw[vec, cA, line width=0.8pt] (\xx,-0.35) -- ++(0,-0.5);
  \foreach \xx in {-0.35, 0, 0.35}
    \draw[vec, cA, line width=0.8pt] (\xx,-2.75) -- ++(0,-0.75);
  \node[lbl, cA, anchor=west] at (0.62,-0.62) {$\vind$};
  \node[lbl, cA, anchor=west] at (0.8,-3.1) {$w$};
\end{tikzpicture}

\caption{The actuator disk in hover. Air at rest far upstream (station 0) is
drawn into a contracting streamtube, crosses the disk (stations 1 and 2) at the
induced velocity $\vind$ with a jump $\Delta p$ in static pressure, and leaves
in a wake of speed $w$ and area $A_\infty$ at ambient pressure (station 3). The
momentum theorem is applied to the dashed control volume. The tube is drawn
from continuity, $r^2u = R^2\vind$, with the axial velocity $u$ of the
vortex-cylinder solution on the axis; it contracts to $A_\infty = A/2$, as
\cref{thm:momentum} requires.}
\label{fig:streamtube}
\end{figure}

\subsection{The three conservation laws}

\paragraph{Mass.} The streamtube carries a constant mass flow. Evaluating at
the disk,
\begin{equation}
  \dot m_{\mathrm{air}} = \rho A \vind = \rho A_\infty w .
  \label{eq:continuity}
\end{equation}

\paragraph{Momentum.} Applying the momentum theorem to a control volume
enclosing the streamtube between a station far upstream, where the fluid is at
rest and the pressure is $p_\infty$, and a station far downstream, where the
pressure has returned to $p_\infty$ and the velocity is $w$, the only external
axial force is that exerted by the disk on the fluid, equal and opposite to the
thrust:
\begin{equation}
  T = \dot m_{\mathrm{air}}\,(w - 0) = \rho A \vind w .
  \label{eq:momentum}
\end{equation}

\paragraph{Energy.} Bernoulli's equation holds along streamlines on each side
of the disk separately, but not across it, because the disk does work on the
fluid. Upstream, from the far field to the disk face,
\begin{equation}
  p_\infty = p_1 + \tfrac12 \rho \vind^2 ,
  \label{eq:bern1}
\end{equation}
and downstream, from the disk face to the far wake,
\begin{equation}
  p_2 + \tfrac12 \rho \vind^2 = p_\infty + \tfrac12 \rho w^2 .
  \label{eq:bern2}
\end{equation}
Subtracting \cref{eq:bern1} from \cref{eq:bern2} eliminates $p_\infty$ and
$\vind$:
\begin{equation}
  \Delta p \coloneqq p_2 - p_1 = \tfrac12 \rho w^2 ,
  \qquad\text{whence}\qquad
  T = A\,\Delta p = \tfrac12 \rho A w^2 .
  \label{eq:pressurejump}
\end{equation}

\subsection{The central result}

\begin{theorem}[Froude's theorem and the hover power law]
\label{thm:momentum}
Under \cref{mod:disk} the far-wake velocity is exactly twice the velocity
induced at the disk,
\begin{equation}
  w = 2\vind ,
  \label{eq:froude}
\end{equation}
the thrust is
\begin{equation}
  T = 2\rho A \vind^2 ,
  \qquad\text{equivalently}\qquad
  \vind = \sqrt{\frac{T}{2\rho A}} ,
  \label{eq:vi}
\end{equation}
the far wake contracts to half the disk area, $A_\infty = A/2$, and the power
delivered to the fluid is
\begin{equation}
  \Pideal = T\vind = \frac{T^{3/2}}{\sqrt{2\rho A}} .
  \label{eq:pideal}
\end{equation}
\end{theorem}

\begin{proof}
Equate the two expressions for thrust, \cref{eq:momentum,eq:pressurejump}:
\[
  \rho A \vind w = \tfrac12 \rho A w^2 .
\]
Since $T>0$ forces $w>0$ and $A>0$, division by $\rho A w$ is legitimate and
gives $\vind = w/2$, which is \cref{eq:froude}. Substituting $w = 2\vind$ into
\cref{eq:momentum} yields $T = 2\rho A \vind^2$, and solving for $\vind$ gives
\cref{eq:vi}; the positive root is selected because the wake is directed
downward for positive thrust. Wake contraction follows from
\cref{eq:continuity}: $A_\infty = A\vind/w = A/2$.

For the power, the rate of work done by the disk on the fluid equals the flux
of kinetic energy out of the control volume, the fluid entering at rest:
\[
  P = \tfrac12 \dot m_{\mathrm{air}} w^2
    = \tfrac12 (\rho A \vind)(2\vind)^2
    = 2\rho A \vind^3
    = (2\rho A \vind^2)\,\vind = T\vind ,
\]
so $P = T\vind$, which is also what one obtains directly as force times the
velocity of the disk relative to the fluid it acts on. Substituting
\cref{eq:vi} gives $\Pideal = T\sqrt{T/(2\rho A)} = T^{3/2}(2\rho A)^{-1/2}$.
\end{proof}

\begin{remark}
The factor of two in \cref{eq:froude} is the reason the disk sees only half the
dynamic pressure of the wake, and is the origin of the $3/2$ exponent that
drives the entire feasibility analysis. Half the velocity increment happens
upstream of the disk and half downstream; \cref{fig:axisprofile} shows both
halves.
\end{remark}

\begin{figure}[tbp]
\centering
% Axial velocity and static pressure along the rotor axis, from Bernoulli on
% each side of the disk with the vortex-cylinder axial velocity
% u/v_i = 1 + xi/sqrt(1+xi^2). Pressure in units of rho v_i^2.
\begin{tikzpicture}[declare function={u(\s) = 1 + \s/sqrt(1+\s*\s);}]
\begin{groupplot}[group style={group size=1 by 2, vertical sep=0.45cm,
                  xlabels at=edge bottom, xticklabels at=edge bottom},
    paperplot, height=4.1cm, xmin=-4, xmax=5, samples=160,
    xlabel={$x/R$, distance along the axis, positive downstream}]
  \nextgroupplot[ylabel={$u/\vind$}, ymin=0, ymax=2.25, ytick={0,1,2}]
    \addplot[cA, domain=-4:5] {u(x)};
    \addplot[construct, domain=-4:5] {2};
    \addplot[construct, domain=-4:5] {1};
    \draw[cG, dashed] (axis cs:0,0) -- (axis cs:0,2.25);
    \node[lbl, anchor=south west] at (axis cs:0.1,0.05) {disk};
    \node[lbl, anchor=north east, fill=white, inner sep=1pt] at (axis cs:4.9,1.88) {$w = 2\vind$};
    \draw[dim] (axis cs:-2.2,0) -- node[lbl, left, fill=white, inner sep=1pt]
        {half upstream} (axis cs:-2.2,1);
    \draw[dim] (axis cs:2.2,1) -- node[lbl, right, fill=white, inner sep=1pt]
        {half downstream} (axis cs:2.2,2);
  \nextgroupplot[ylabel={$(p-p_\infty)/\rho\vind^2$}, ymin=-2.25, ymax=1.9,
                 ytick={-2,-0.5,0,1.5}]
    \addplot[cA, domain=-4:0] {-0.5*u(x)^2};
    \addplot[cA, domain=0:5] {2 - 0.5*u(x)^2};
    \addplot[construct, domain=-4:5] {0};
    \draw[cD, line width=1.1pt, -{Stealth[length=2mm]}] (axis cs:0,-0.5) -- (axis cs:0,1.5);
    \node[lbl, cD, anchor=east, fill=white, inner sep=1pt] at (axis cs:-0.15,0.75)
        {$\Delta p = 2\rho\vind^2 = \tfrac12\rho w^2$};
    \node[lbl, anchor=north, fill=white, inner sep=1pt] at (axis cs:-2.4,-1.0)
        {suction ahead of the disk};
    \node[lbl, anchor=north, fill=white, inner sep=1pt] at (axis cs:3.0,-1.0)
        {recovers to $p_\infty$ in the wake};
\end{groupplot}
\end{tikzpicture}

\caption{Axial velocity and static pressure along the rotor axis, from
Bernoulli's equation applied separately on each side of the disk with the
vortex-cylinder velocity $u/\vind = 1 + \xi/\sqrt{1+\xi^2}$, $\xi = x/R$. The
velocity is continuous through the disk and rises from $0$ to $\vind$
upstream and from $\vind$ to $w = 2\vind$ downstream, the two halves of
\cref{eq:froude}. The pressure falls ahead of the disk, jumps by
$\Delta p = \tfrac12\rho w^2$ across it, \cref{eq:pressurejump}, and recovers
to $p_\infty$ in the wake.}
\label{fig:axisprofile}
\end{figure}

\begin{corollary}[Hover]
\label{cor:hover}
In hover, $T = mg$, so
\begin{equation}
  \vind = \sqrt{\frac{mg}{2\rho A}},
  \qquad
  \Pideal = \frac{(mg)^{3/2}}{\sqrt{2\rho A}} .
  \label{eq:hoverpower}
\end{equation}
At fixed disk area, $\Pideal \propto m^{3/2}$.
\end{corollary}

\subsection{Disk loading}

\begin{definition}[Disk loading]
$\mathrm{DL} \coloneqq T/A$, with units of pressure.
\end{definition}

Rewriting \cref{eq:vi,eq:pideal} in terms of disk loading isolates the
physically meaningful group:
\begin{equation}
  \vind = \sqrt{\frac{\mathrm{DL}}{2\rho}},
  \qquad
  \frac{\Pideal}{T} = \sqrt{\frac{\mathrm{DL}}{2\rho}} ,
  \qquad
  \frac{\Pideal}{m} = g\sqrt{\frac{\mathrm{DL}}{2\rho}} .
  \label{eq:dl}
\end{equation}

\begin{observation}
\label{obs:dl}
Power per unit mass in hover depends on the disk loading alone, not on mass and
area separately. A hovering machine is therefore not primarily asking for
powerful motors; it is asking for low disk loading. The quantity
$T/\Pideal = \sqrt{2\rho/\mathrm{DL}}$, the power loading, is the correct
figure of merit for a hovering vehicle and improves as disk loading falls.
\end{observation}

\Cref{fig:powerloading} places the design of \cref{sec:results} on this law.

\begin{figure}[tbp]
\centering
% Power loading against disk loading, T/P = FM sqrt(2 rho / DL), eq:dl.
\providecommand{\aThover}{4.23}
\providecommand{\aVstall}{24.71}
\providecommand{\aVdesign}{26.69}
\providecommand{\aPstall}{17.04}
\providecommand{\aPdesign}{17.56}
\providecommand{\aPind}{14.87}
\providecommand{\aCtMax}{0.14}
\providecommand{\aCtDesign}{0.12}
\providecommand{\aSigma}{0.212}
\providecommand{\aRho}{1.225}
\providecommand{\aDL}{22.2}
\providecommand{\aWake}{6.02}
\providecommand{\aKappa}{1.168}

\begin{tikzpicture}
\begin{axis}[paperplot, xmode=log, ymode=log, xmin=5, xmax=2000,
    ymin=0.02, ymax=1.3, samples=40, domain=5:2000,
    xlabel={disk loading $\mathrm{DL} = T/A$ [\si{\pascal}]},
    ylabel={power loading $T/P$ [\si{\newton\per\watt}]},
    legend pos=south west]
  % wake faster than 5 m/s: loose paper moves, requirement (R3)
  \pgfmathsetmacro{\DLfive}{2*\aRho*2.5*2.5}
  \fill[cD!7] (axis cs:\DLfive,0.02) rectangle (axis cs:2000,1.3);
  \draw[cD, dashed, line width=0.6pt] (axis cs:\DLfive,0.02) -- (axis cs:\DLfive,1.3);
  \node[lbl, cD, anchor=north west] at (axis cs:\DLfive*1.05,1.25)
      {wake $w = 2\vind > \SI{5}{\metre\per\second}$};
  \addplot[cA] {sqrt(2*\aRho/x)};
  \addlegendentry{$\FM = 1$, ideal disk}
  \addplot[cA!70, dashed] {0.75*sqrt(2*\aRho/x)};
  \addlegendentry{$\FM = 0.75$}
  \addplot[cA!45, dotted, line width=1.2pt] {0.5*sqrt(2*\aRho/x)};
  \addlegendentry{$\FM = 0.5$}
  % the design
  \pgfmathsetmacro{\TPd}{\rFM*sqrt(2*\aRho/\aDL)}
  \addplot[only marks, mark=*, mark size=2.2pt, cA, forget plot]
      coordinates {(\aDL,\TPd)};
  \node[lbl, anchor=north west, align=left, fill=white, inner sep=1.5pt]
      (dl) at (axis cs:40,0.62)
      {design: $\mathrm{DL}=\SI{\aDL}{\pascal}$,\\
       $\FM=\rFM$, $w=\SI{\aWake}{\metre\per\second}$};
  \draw[cG, line width=0.4pt] (dl.south west) -- (axis cs:\aDL*1.06,\TPd*1.04);
  % slope triangle, -1/2
  \draw[cG, line width=0.5pt] (axis cs:400,0.16) -- (axis cs:1600,0.16)
      -- (axis cs:400,0.32) -- cycle;
  \node[lbl, cG, anchor=north] at (axis cs:800,0.155) {slope $-\tfrac12$};
\end{axis}
\end{tikzpicture}

\caption{Power loading against disk loading, $T/P = \FM\sqrt{2\rho/\mathrm{DL}}$,
from \cref{eq:dl} and the definition of $\FM$ in \cref{sec:blade}. On
logarithmic axes every figure of merit gives a line of slope $-\tfrac12$:
halving the disk loading buys a factor $\sqrt2$ in thrust per watt, whatever
the rotor. The shaded region is where the wake exceeds
\SI{5}{\metre\per\second}, the speed at which loose paper begins to move
(requirement (R3) of \cref{sec:results}); the design sits just inside it, at
$\mathrm{DL} = \SI{\aDL}{\pascal}$ and $w = \SI{\rWake}{\metre\per\second}$.}
\label{fig:powerloading}
\end{figure}

This single observation determines the shape of every machine in
\cref{sec:results}: large, slowly turning rotors, because
\cref{eq:dl} says that is the only lever with any authority, and
\cref{sec:acoustics} will show that the same choice is what makes the machine
quiet.

\subsection{Axial flight and the validity of \texorpdfstring{\cref{thm:momentum}}{Theorem 3.2}}

For axial motion at climb velocity $V_c$ relative to the air, the same control
volume argument with a non-zero inflow gives
\begin{equation}
  T = 2\rho A (V_c + \vind)\vind
  \quad\Longrightarrow\quad
  \vind = -\frac{V_c}{2} + \sqrt{\left(\frac{V_c}{2}\right)^2
          + \frac{T}{2\rho A}} ,
  \label{eq:axial}
\end{equation}
which reduces to \cref{eq:vi} at $V_c=0$ and is implemented in that form. The
power is then $P = T(V_c + \vind)$.

\Cref{eq:axial} is valid in climb and in the windmill-brake state, but fails in
descent for $-2 \lesssim V_c/v_{i,h} \lesssim 0$, where $v_{i,h}$ is the hover
induced velocity. This is the vortex ring state, where
the streamtube assumption breaks down and the flow is unsteady and
recirculating~\cite[\S2.13]{leishman2006}. The machines considered here descend
slowly and the engine does not enter that regime, but it is a real limitation
of the model rather than a technicality: a fast descent to land is precisely
the manoeuvre it cannot predict.

Two further caveats are worth recording. Momentum theory gives a lower bound on
the power: it is the power required to produce the thrust by the most efficient
possible redistribution of momentum, and any real rotor does worse.
\Cref{sec:blade} quantifies by how much. Second, the theory says nothing about
how the loading is produced, so it cannot predict stall, compressibility or
Reynolds effects; those enter through the blade element layer.
